\section{ماژول huffman\_coder (Huffman Coder)}
\label{sec:module_huffman_coder}

\subsection{معرفی و انگیزه}
این مکانیزم بر مبنای مبانی نظری پیشرفته در ادبیات توجه و فشرده‌سازی \cite{ref_huffmancoder_2024} توسعه یافته است.

کوانتایش وزن‌ها و فعال‌سازی‌ها روشی کلیدی در فشرده‌سازی مدل‌های یادگیری عمیق است که مقادیر ممیز شناور ۳۲ یا ۱۶ بیتی را به کدهای گسسته با طول بیت پایین (مانند ۴ بیتی یا ۸ بیتی) تبدیل می‌کند. با این حال، حتی پس از کوانتایش یکنواخت، توزیع فراوانی نمادها یا سطوح کوانتایز در شبکه‌های عصبی هرگز کاملاً یکنواخت نبوده و مقادیر نزدیک به صفر فراوانی بسیار بالاتری دارند. بنابراین، ذخیره‌سازی نمادها با طول بیت ثابت، اصل آنتروپی اطلاعات شانون را نادیده گرفته و به هدررفت ظرفیت ذخیره‌سازی می‌انجامد.

موتور فشرده‌سازی \lr{Huffman Coder (Entropy Compression Engine)} به عنوان یک الگوریتم کدگذاری آنتروپی بدون اتلاف (\lr{Lossless Entropy Coding}) برای فشرده‌سازی لایه‌های کوانتایز شده و همچنین حافظه نهان کلید و مقدار (\lr{KV Cache}) طراحی شده است. این ماژول با ساخت درخت پیشوندی بهینه هافمن متناسب با توزیع آماری وزن‌ها، طول کدهای باینری را معکوس با احتمال وقوع نمادها تنظیم نموده و امکان نزدیک شدن نرخ ذخیره‌سازی به حد آنتروپی شانون را فراهم می‌سازد.

\subsection{مبانی نظری و فرمول‌بندی ریاضی}
فرض کنید تانسور وزن یا مقادیر کوانتایز شده یک لایه به صورت مجموعه‌ای از نمادها $\mathcal{S} = \{s_1, s_2, \dots, s_K\}$ با توزیع احتمال تجربی $p(s_k) = \frac{n_k}{N}$ باشد، که در آن $N$ تعداد کل المان‌ها و $n_k$ فراوانی نماد $s_k$ است. آنتروپی شانون مجموعه داده برابر است با:
\begin{equation}
H(\mathcal{S}) = -\sum_{k=1}^K p(s_k) \log_2 p(s_k) \quad \text{(bits/symbol)}
\end{equation}
الگوریتم کدگذاری هافمن با تخصیص طول کد متغیر $\ell(s_k) \in \mathbb{N}$ به هر نماد، مسئله کمینه‌سازی طول مورد انتظار کلمات کد را حل می‌کند:
\begin{equation}
\min_{\ell} \bar{L} = \sum_{k=1}^K p(s_k) \ell(s_k) \quad \text{subject to} \quad \sum_{k=1}^K 2^{-\ell(s_k)} \le 1 \quad \text{(Kraft-McMillan Inequality)}
\end{equation}
با حل بازگشتی از طریق صف اولویت‌دار (\lr{Min-Heap Priority Queue})، درخت دوتایی پیشوندی هافمن تشکیل شده و طول کد ارائه‌شده ارضای رابطه $H(\mathcal{S}) \le \bar{L} < H(\mathcal{S}) + 1$ را تضمین می‌کند. در فاز بازیابی و دیکودینگ برداری، جریان بیتی فشرده‌شده با پیمایش پویای درخت بازسازی می‌گردد. در فاز بازسازی لایه‌های شبکه و آموزش کم‌مرتبه، مدل از لایه‌های تثبیت‌شده به همراه به‌روزرسانی آداپتورهای گرادیانی به فرم زیر استفاده می‌کند:
\begin{equation}
\mathbf{y} = \mathbf{x} \mathbf{W}_{\text{quant}}^\top + \mathbf{b} + \frac{\alpha}{r_{\text{lora}}} (\mathbf{x} \mathbf{A}^\top) \mathbf{B}^\top
\end{equation}

\paragraph{تحلیل طیفی و نامساوی‌های پایداری خطا}
برای بررسی میزان دقت تقریب ماتریس وزن یا نگاشت توجه در ماژول \lr{huffman\_coder}، تفاضل ماتریس فاکتورگیری‌شده از نگاشت متراکم اصلی بر اساس نامساوی هافدینگ-فون‌نیومان ارزیابی می‌گردد. فرض کنید $\mathbf{W} \in \mathbb{R}^{d_{\text{in}} \times d_{\text{out}}}$ ماتریس وزن هدف باشد که توسط تقریب کم‌مرتبه یا پراکنده با رتبه موثر $r$ جایگزین شده است. کران بالای خطای بازسازی بر مبنای نرم فروبنیوس به صورت زیر تصریح می‌شود:
\begin{equation}
\label{eq:spectral_bound_huffman_coder}
\| \mathbf{W} - \mathbf{W}_{r} \|_F = \sqrt{\sum_{i=r+1}^{\min(d_{\text{in}}, d_{\text{out}})} \sigma_i^2(\mathbf{W})} \le \epsilon \| \mathbf{W} \|_F
\end{equation}
که در آن $\sigma_i(\mathbf{W})$ بیانگر $i$-امین مقدار منفرد مرتب‌شده ماتریس به صورت نزولی است. نسبت حفظ انرژی طیفی $\eta_r = \frac{\sum_{i=1}^r \sigma_i^2}{\sum_{i=1}^{d} \sigma_i^2}$ تضمین می‌کند که با انتخاب آستانه $\tau \ge 0.90$، بیش از نود درصد از پراکندگی واریانس داده‌ها حفظ شده و افت سیگنال در انتشار گرادیان رو به عقب به حداقل ممکن کاهش می‌یابد.

\paragraph{تحلیل مجانبی و پیچیدگی اشتراک حافظه سخت‌افزاری}
از منظر پیچیدگی محاسباتی و انتقال داده‌ها در سلسومراتب حافظه (حافظه سراسری \lr{HBM} به حافظه محلی \lr{SRAM})، عملیات ضرب ماتریسی در این مکانیزم نیازمند تخصیص بافر موقت با ابعاد فاکتورگیری‌شده است. برای پردازش دنباله‌ای به طول $N$ با بعد مخفی $D$، بار محاسباتی هسته از فرم درجه دو به رابطه زیر تقلیل می‌یابد:
\begin{equation}
\label{eq:complexity_huffman_coder}
\mathcal{C}_{\text{huffman\_coder}} = \mathcal{O}\left( B \cdot N \cdot D \cdot r + B \cdot H \cdot N \cdot d_k \right) \ll \mathcal{O}\left( B \cdot H \cdot N^2 \cdot d_k \right)
\end{equation}
این بهبود مرتبه پیچیدگی، با کاهش چشمگیر حجم تبادل تانسورها در گذرگاه‌های ارتباطی کارت گرافیک، مانع از اشباع پهنای باند حافظه (پدیده \lr{Memory-Bound}) در فرآیند استنتاج طولانی‌مدت می‌گردد.

\subsection{معماری و پیاده‌سازی}
پیاده‌سازی این سیستم در فایل \lr{Codes/framework/compression\_engine/huffman\_coder.py} انجام یافته و در رجیستری فشرده‌سازی با نام‌های \lr{huffman} و \lr{huffman\_coder} به ثبت رسیده است. ویژگی‌های معماری موتور شامل موارد زیر است:
\begin{itemize}
    \item \textbf{گره و درخت هافمن (\lr{HuffmanNode} و \lr{HuffmanCoder}):} ماژول تخصصی جهت محاسبه برداری فراوانی نمادها، ساخت سریع درخت پیشوندی در حافظه و بسته‌بندی بیت‌ها در آرایه‌های فشرده \lr{numpy.uint8}.
    \item \textbf{کوانتایش پیوسته وزن‌ها (\lr{Uniform Weight Quantization}):} نگاشت متقارن مقادیر ممیز شناور به کدهای صحیح با تعداد بیت‌های قابل تنظیم ($b \in \{4, 8\}$).
    \item \textbf{تثبیت وزن پایه و حفظ گرادیان \lr{LoRA}:} تابع \lr{\_freeze\_base\_only} تضمین می‌کند که پارامترهای اصلی کوانتایز شده تثبیت گردیده و تنها پارامترهای تنظیمی \lr{LoRA} در طول بازیابی گرادیانی آموزش‌پذیر باقی بمانند.
    \item \textbf{حسابرسی دقیق بایت‌ها (\lr{Exact Byte Accounting}):} محاسبه نرخ فشرده‌سازی فیزیکی بر اساس تعداد بایت‌های واقعی فشرده‌شده و جداول کدهای ذخیره‌شده.
\end{itemize}

\subsection{نتایج تجربی}
ارزیابی آزمایشگاهی بنچ‌مارک فشرده‌ساز \lr{huffman\_coder} بر روی مدل‌های زبانی در جدول \ref{tab:benchmark_huffman_coder} گزارش شده است. این نتایج به صورت مستقیم از داده‌های اعتبارسنجی پروژه برداشت گردیده است:

\begin{table}[H]
\centering
\small
\caption{نتایج بنچ‌مارک تجربی فشرده‌ساز \lr{huffman\_coder} بر روی مدل‌های پایه}
\label{tab:benchmark_huffman_coder}
\begin{tabular}{lccccc}
\toprule
\textbf{مدل پایه} \lr{(Base Model)} & \textbf{دقت} \lr{(Accuracy)} & \textbf{شاخص} \lr{(F1-Score)} & \textbf{حافظه} \lr{(Peak VRAM)} & \textbf{توان} \lr{(Throughput)} & \textbf{وضعیت} \\
\midrule
\lr{DistilBERT} & \lr{90.60\%} & \lr{0.9095} & \lr{308.84 MB} & \lr{158,785.10 tok/s} & \lr{PASS} \\
\lr{RoBERTa} & \lr{51.38\%} & \lr{0.6768} & \lr{528.86 MB} & \lr{86,035.52 tok/s} & \lr{PASS\_LOW} \\
\lr{GPT-2} & \lr{50.92\%} & \lr{0.6748} & \lr{634.91 MB} & \lr{74,202.00 tok/s} & \lr{PASS\_LOW} \\
\lr{Qwen} & \lr{0.00\%} & \lr{---} & \lr{---} & \lr{---} & \lr{SKIP} \\
\bottomrule
\end{tabular}
\end{table}

همان‌گونه که در جدول مشخص است، فشرده‌ساز \lr{huffman\_coder} بر روی مدل \lr{DistilBERT} پایداری کم‌نظیری داشته و دقت \lr{90.60\%} همراه با شاخص F1 معادل \lr{0.9095} و مصرف حافظه ناچیز \lr{308.84 MB} را با وضعیت موفقیت‌آمیز \lr{PASS} به ثبت رسانده است. در مدل‌های \lr{RoBERTa} و \lr{GPT-2}، کوانتایش بدون کالیبراسیون دقیق لایه‌نرمال افت دقت را به دنبال داشته و به ترتیب دقت‌های \lr{51.38\%} و \lr{50.92\%} را حاصل نموده‌اند که در دروازه کیفیت تحت عنوان \lr{PASS\_LOW} علامت‌گذاری شده‌اند. ارزیابی مدل \lr{Qwen} به علت ناسازگاری رجیستری در پیکربندی و اجرای ناموفق به عنوان \lr{SKIP} گزارش شده است.

\subsection{محدودیت‌ها و تحلیل شکست}
تحلیل عملکردی ماژول نقاط ضعف زیر را نمایان می‌سازد:
\begin{itemize}
    \item \textbf{افت دقت ناشی از کوانتایش تکمیلی در مدل‌های عمیق (\lr{RoBERTa} و \lr{GPT-2}):} کدگذاری هافمن به خودی خود بدون اتلاف است، اما اتلاف دقت از مرحله پیش‌نیاز کوانتایش گسسته ناشی می‌شود. در مدل‌های با لایه‌های متوالی متعدد، تجمع خطای کوانتایش در تمام لایه‌ها موجب افت دقت به سطح پایه تصادفی شده و وضعیت \lr{PASS\_LOW} را ایجاد می‌کند.
    \item \textbf{ناهمگونی تاخیر پردازشی در دیکودینگ ناهمگن (\lr{Variable Bit-Length Decoding}):} به علت طول متغیر کدهای باینری، دیکودینگ موازی در سطح رشته‌های موازی (\lr{SIMD/Warp}) در پردازنده‌های گرافیکی دچار انحراف انشعاب (\lr{Warp Divergence}) شده و نیازمند پیاده‌سازی کرنل‌های اختصاصی سخت‌افزاری است.
\end{itemize}
